% ============================================
% SECÇÃO 3.3 - MODELAÇÃO COMPORTAMENTAL
% Sistema de Gestão Integrada - Clube Amigos de Polvoreira
% ============================================

\subsection{Modelação comportamental}

Foram construídos vários diagramas que procuravam representar o comportamento do sistema. 
O primeiro deles foi um diagrama de máquina de estados, que descreve os estados em que uma 
quota se pode encontrar, e que ações a fazem mudar de estado. O motivo para a construção deste 
diagrama foi a dificuldade de explicar estes estados e as transições entre eles diretamente a partir dos 
requisitos, uma vez que a informação necessária para o fazer se encontra espalhada por vários 
requisitos, dificultando que um membro da equipa de desenvolvimento compreenda esta máquina de 
estados sem uma representação concisa da mesma.

% Diagrama de máquina de estados - Quota
\begin{figure}[H]
    \centering
    \includegraphics[width=0.8\textwidth]{imagens/maquina_estados_quota.png}
    \caption{Diagrama de máquina de estados dos estados de uma quota.}
    \label{fig:maquina_estados_quota}
\end{figure}

Uma quota é criada no início de cada época desportiva com o estado ``Pendente'', aguardando o seu pagamento. 
Quando o encarregado de educação ou atleta efetua o pagamento, a quota transita para o estado ``Paga''. 
Caso o prazo de pagamento expire sem que o pagamento seja efetuado, a quota transita para o estado ``Em Atraso''. 
A partir deste estado, o pagamento ainda pode ser realizado, transitando para ``Paga'', ou pode ser aplicada 
uma suspensão ao atleta caso o atraso persista. Uma quota no estado ``Paga'' pode ainda ser reembolsada 
em casos excecionais, transitando para o estado ``Reembolsada''.

De forma análoga, foi construído um diagrama de máquina de estados para representar os possíveis 
estados de uma convocatória:

% Diagrama de máquina de estados - Convocatória
\begin{figure}[H]
    \centering
    \includegraphics[width=0.8\textwidth]{imagens/maquina_estados_convocatoria.png}
    \caption{Diagrama de máquina de estados dos estados de uma convocatória.}
    \label{fig:maquina_estados_convocatoria}
\end{figure}

Uma convocatória é criada pelo treinador com o estado ``Rascunho'', permitindo a sua edição. 
Quando o treinador finaliza a seleção de atletas, a convocatória é publicada, transitando para o estado ``Publicada''. 
Neste estado, os encarregados de educação podem confirmar ou recusar a presença dos seus educandos. 
Quando o jogo ou treino associado termina, a convocatória transita automaticamente para o estado ``Finalizada''. 
Uma convocatória pode ainda ser cancelada a partir dos estados ``Rascunho'' ou ``Publicada''.

% ============================================
% DIAGRAMA DE CASOS DE USO GERAL
% ============================================

De seguida, foi construído um diagrama de casos de uso, que apresenta todas as operações suportadas 
pelo sistema, e que atores as podem realizar. Devido ao elevado número de casos de uso, dividiu-se o 
diagrama de acordo com os vários subsistemas. Apresenta-se, abaixo, um diagrama geral, onde se 
encontram representados os vários atores e os subsistemas com que interagem.

\begin{figure}[H]
    \centering
    \includegraphics[width=1.0\textwidth]{imagens/casos_uso_geral.png}
    \caption{Diagrama de casos de uso geral.}
    \label{fig:casos_uso_geral}
\end{figure}

No diagrama, observam-se os vários atores que interagem com o sistema:
\begin{itemize}
    \item \textbf{Utilizador} -- ator genérico que representa operações comuns a todos os tipos de utilizador
    \item \textbf{Atleta} -- praticante de uma ou mais modalidades do clube
    \item \textbf{Treinador} -- responsável por equipas e gestão de treinos/jogos
    \item \textbf{Encarregado de Educação} -- responsável legal por atletas menores
    \item \textbf{Secretaria} -- funcionário responsável pela gestão administrativa
    \item \textbf{Gerência} -- administrador do sistema com acesso total
\end{itemize}

Apresentam-se, abaixo, os diagramas de casos de uso correspondentes aos vários subsistemas:

% ============================================
% DIAGRAMA DE CASOS DE USO - GESTÃO DE UTILIZADORES
% ============================================

\begin{figure}[H]
    \centering
    \includegraphics[width=0.9\textwidth]{imagens/casos_uso_utilizadores.png}
    \caption{Diagrama de casos de uso do subsistema de Gestão de Utilizadores.}
    \label{fig:casos_uso_utilizadores}
\end{figure}

% ============================================
% DIAGRAMA DE CASOS DE USO - GESTÃO DESPORTIVA
% ============================================

\begin{figure}[H]
    \centering
    \includegraphics[width=0.9\textwidth]{imagens/casos_uso_desportiva.png}
    \caption{Diagrama de casos de uso do subsistema de Gestão Desportiva.}
    \label{fig:casos_uso_desportiva}
\end{figure}

% ============================================
% DIAGRAMA DE CASOS DE USO - GESTÃO CLÍNICA
% ============================================

\begin{figure}[H]
    \centering
    \includegraphics[width=0.9\textwidth]{imagens/casos_uso_clinica.png}
    \caption{Diagrama de casos de uso do subsistema de Gestão Clínica.}
    \label{fig:casos_uso_clinica}
\end{figure}

% ============================================
% DIAGRAMA DE CASOS DE USO - GESTÃO FINANCEIRA
% ============================================

\begin{figure}[H]
    \centering
    \includegraphics[width=0.9\textwidth]{imagens/casos_uso_financeira.png}
    \caption{Diagrama de casos de uso do subsistema de Gestão Financeira.}
    \label{fig:casos_uso_financeira}
\end{figure}

% ============================================
% DESCRIÇÃO DOS CASOS DE USO
% ============================================

\subsubsection{Descrição dos casos de uso}

Seguem-se as descrições detalhadas dos casos de uso identificados, organizadas por subsistema.

% --- GESTÃO DE UTILIZADORES ---

\paragraph{Subsistema de Gestão de Utilizadores}

\begin{table}[H]
\centering
\begin{tabularx}{\textwidth}{|l|X|}
\hline
\textbf{Caso de uso} & Registar-se \\
\hline
\textbf{Ator} & Utilizador \\
\hline
\textbf{Descrição} & Utilizador cria uma conta para se poder autenticar e realizar outras operações. \\
\hline
\textbf{Pré-condição} & Utilizador não tem sessão iniciada. \\
\hline
\textbf{Pós-condição} & É registada uma conta com os dados providenciados pelo utilizador, pendente de aprovação. \\
\hline
\textbf{Fluxo normal} & 
1. Utilizador solicita o seu registo. \newline
2. Utilizador indica o seu nome, email, palavra-passe e tipo de conta pretendido. \newline
3. Sistema verifica que não existe um utilizador com o email providenciado. \newline
4. Sistema regista os dados do utilizador numa nova conta com estado ``Pendente''. \newline
5. Sistema notifica a secretaria sobre o novo pedido de registo. \newline
6. Sistema informa utilizador de sucesso na operação de registo. \\
\hline
\textbf{Fluxo de exceção 1} & 
[existe um utilizador com o email providenciado] (passo 3) \newline
3.1. Sistema informa o utilizador que o email já corresponde a um utilizador registado. \newline
3.2. Sistema cancela a operação de registo. \\
\hline
\end{tabularx}
\caption{Caso de uso ``Registar-se''.}
\label{tab:uc_registar}
\end{table}

\begin{table}[H]
\centering
\begin{tabularx}{\textwidth}{|l|X|}
\hline
\textbf{Caso de uso} & Iniciar Sessão \\
\hline
\textbf{Ator} & Utilizador \\
\hline
\textbf{Descrição} & Utilizador autentica-se para poder realizar outras operações. \\
\hline
\textbf{Pré-condição} & Utilizador não tem sessão iniciada e a sua conta está aprovada. \\
\hline
\textbf{Pós-condição} & Utilizador fica com sessão iniciada no sistema. \\
\hline
\textbf{Fluxo normal} & 
1. Utilizador providencia email e palavra-passe. \newline
2. Sistema verifica que existe um utilizador com as credenciais providenciadas e conta ativa. \newline
3. Sistema inicia uma sessão para o utilizador. \newline
4. Sistema redireciona o utilizador para a página principal correspondente ao seu tipo. \\
\hline
\textbf{Fluxo de exceção 1} & 
[credenciais inválidas ou conta não ativa] (passo 2) \newline
2.1. Sistema informa que a autenticação não teve sucesso. \newline
2.2. Sistema cancela a operação de autenticação. \\
\hline
\end{tabularx}
\caption{Caso de uso ``Iniciar Sessão''.}
\label{tab:uc_iniciar_sessao}
\end{table}

\begin{table}[H]
\centering
\begin{tabularx}{\textwidth}{|l|X|}
\hline
\textbf{Caso de uso} & Aprovar Registo \\
\hline
\textbf{Ator} & Secretaria \\
\hline
\textbf{Descrição} & Secretaria aprova ou rejeita um pedido de registo de utilizador. \\
\hline
\textbf{Pré-condição} & Secretaria tem sessão iniciada e existem pedidos de registo pendentes. \\
\hline
\textbf{Pós-condição} & O pedido de registo é aprovado ou rejeitado, e o utilizador é notificado. \\
\hline
\textbf{Fluxo normal} & 
1. Secretaria consulta a lista de pedidos de registo pendentes. \newline
2. Secretaria seleciona um pedido para análise. \newline
3. Sistema apresenta os dados do pedido. \newline
4. Secretaria aprova o pedido. \newline
5. Sistema atualiza o estado da conta para ``Ativo''. \newline
6. Sistema envia notificação ao utilizador sobre a aprovação. \\
\hline
\textbf{Fluxo alternativo 1} & 
[secretaria rejeita o pedido] (passo 4) \newline
4.1. Secretaria indica o motivo da rejeição. \newline
4.2. Sistema atualiza o estado da conta para ``Rejeitado''. \newline
4.3. Sistema envia notificação ao utilizador sobre a rejeição. \\
\hline
\end{tabularx}
\caption{Caso de uso ``Aprovar Registo''.}
\label{tab:uc_aprovar_registo}
\end{table}

\begin{table}[H]
\centering
\begin{tabularx}{\textwidth}{|l|X|}
\hline
\textbf{Caso de uso} & Gerir Perfil \\
\hline
\textbf{Ator} & Utilizador \\
\hline
\textbf{Descrição} & Utilizador consulta e edita os seus dados pessoais. \\
\hline
\textbf{Pré-condição} & Utilizador tem sessão iniciada. \\
\hline
\textbf{Pós-condição} & Os dados do perfil são atualizados conforme as alterações efetuadas. \\
\hline
\textbf{Fluxo normal} & 
1. Utilizador acede ao seu perfil. \newline
2. Sistema apresenta os dados atuais do utilizador. \newline
3. Utilizador edita os dados pretendidos. \newline
4. Sistema valida os novos dados. \newline
5. Sistema guarda as alterações. \newline
6. Sistema confirma o sucesso da operação. \\
\hline
\textbf{Fluxo de exceção 1} & 
[dados inválidos] (passo 4) \newline
4.1. Sistema informa quais campos contêm erros. \newline
4.2. Regressa a 3. \\
\hline
\end{tabularx}
\caption{Caso de uso ``Gerir Perfil''.}
\label{tab:uc_gerir_perfil}
\end{table}

% --- GESTÃO DESPORTIVA ---

\paragraph{Subsistema de Gestão Desportiva}

\begin{table}[H]
\centering
\begin{tabularx}{\textwidth}{|l|X|}
\hline
\textbf{Caso de uso} & Criar Convocatória \\
\hline
\textbf{Ator} & Treinador \\
\hline
\textbf{Descrição} & Treinador cria uma convocatória para um jogo ou treino. \\
\hline
\textbf{Pré-condição} & Treinador tem sessão iniciada e está associado a pelo menos uma equipa. \\
\hline
\textbf{Pós-condição} & É criada uma convocatória com os atletas selecionados, em estado ``Rascunho''. \\
\hline
\textbf{Fluxo normal} & 
1. Treinador seleciona a equipa. \newline
2. Sistema apresenta o plantel da equipa. \newline
3. Treinador indica o tipo de evento (jogo/treino), data, hora e local. \newline
4. Treinador seleciona os atletas a convocar. \newline
5. Sistema cria a convocatória em estado ``Rascunho''. \newline
6. Sistema confirma a criação da convocatória. \\
\hline
\end{tabularx}
\caption{Caso de uso ``Criar Convocatória''.}
\label{tab:uc_criar_convocatoria}
\end{table}

\begin{table}[H]
\centering
\begin{tabularx}{\textwidth}{|l|X|}
\hline
\textbf{Caso de uso} & Publicar Convocatória \\
\hline
\textbf{Ator} & Treinador \\
\hline
\textbf{Descrição} & Treinador publica uma convocatória, notificando os atletas e encarregados. \\
\hline
\textbf{Pré-condição} & Treinador tem sessão iniciada e existe uma convocatória em estado ``Rascunho''. \\
\hline
\textbf{Pós-condição} & A convocatória passa ao estado ``Publicada'' e são enviadas notificações. \\
\hline
\textbf{Fluxo normal} & 
1. Treinador seleciona uma convocatória em rascunho. \newline
2. Sistema apresenta os detalhes da convocatória. \newline
3. Treinador confirma a publicação. \newline
4. Sistema atualiza o estado para ``Publicada''. \newline
5. Sistema envia notificações aos atletas convocados e seus encarregados de educação. \\
\hline
\end{tabularx}
\caption{Caso de uso ``Publicar Convocatória''.}
\label{tab:uc_publicar_convocatoria}
\end{table}

\begin{table}[H]
\centering
\begin{tabularx}{\textwidth}{|l|X|}
\hline
\textbf{Caso de uso} & Responder a Convocatória \\
\hline
\textbf{Ator} & Encarregado de Educação \\
\hline
\textbf{Descrição} & Encarregado confirma ou recusa a presença do seu educando numa convocatória. \\
\hline
\textbf{Pré-condição} & Encarregado tem sessão iniciada e existe uma convocatória publicada para o seu educando. \\
\hline
\textbf{Pós-condição} & A resposta do encarregado é registada na convocatória. \\
\hline
\textbf{Fluxo normal} & 
1. Encarregado consulta as convocatórias pendentes de resposta. \newline
2. Sistema apresenta a lista de convocatórias. \newline
3. Encarregado seleciona uma convocatória. \newline
4. Sistema apresenta os detalhes da convocatória. \newline
5. Encarregado confirma a presença do educando. \newline
6. Sistema regista a confirmação. \newline
7. Sistema notifica o treinador sobre a resposta. \\
\hline
\textbf{Fluxo alternativo 1} & 
[encarregado recusa a presença] (passo 5) \newline
5.1. Encarregado indica o motivo da ausência. \newline
5.2. Sistema regista a recusa com a justificação. \newline
5.3. Sistema notifica o treinador sobre a resposta. \\
\hline
\end{tabularx}
\caption{Caso de uso ``Responder a Convocatória''.}
\label{tab:uc_responder_convocatoria}
\end{table}

\begin{table}[H]
\centering
\begin{tabularx}{\textwidth}{|l|X|}
\hline
\textbf{Caso de uso} & Registar Estatísticas de Jogo \\
\hline
\textbf{Ator} & Treinador \\
\hline
\textbf{Descrição} & Treinador regista as estatísticas individuais e coletivas de um jogo. \\
\hline
\textbf{Pré-condição} & Treinador tem sessão iniciada e existe um jogo finalizado sem estatísticas registadas. \\
\hline
\textbf{Pós-condição} & As estatísticas do jogo são registadas no sistema. \\
\hline
\textbf{Fluxo normal} & 
1. Treinador seleciona o jogo. \newline
2. Sistema apresenta a lista de atletas que participaram. \newline
3. Treinador indica o resultado do jogo. \newline
4. Para cada atleta, treinador regista estatísticas individuais (golos, assistências, tempo jogado, etc.). \newline
5. Sistema valida e guarda as estatísticas. \newline
6. Sistema confirma o registo das estatísticas. \\
\hline
\end{tabularx}
\caption{Caso de uso ``Registar Estatísticas de Jogo''.}
\label{tab:uc_registar_estatisticas}
\end{table}

\begin{table}[H]
\centering
\begin{tabularx}{\textwidth}{|l|X|}
\hline
\textbf{Caso de uso} & Consultar Plantel \\
\hline
\textbf{Ator} & Treinador \\
\hline
\textbf{Descrição} & Treinador consulta os atletas que compõem o plantel de uma equipa. \\
\hline
\textbf{Pré-condição} & Treinador tem sessão iniciada e está associado a pelo menos uma equipa. \\
\hline
\textbf{Pós-condição} & São apresentados ao treinador os atletas do plantel com os seus dados. \\
\hline
\textbf{Fluxo normal} & 
1. Treinador seleciona a equipa. \newline
2. Sistema apresenta a lista de atletas do plantel com informações relevantes (posição, idade, atestado médico válido, quotas em dia). \\
\hline
\end{tabularx}
\caption{Caso de uso ``Consultar Plantel''.}
\label{tab:uc_consultar_plantel}
\end{table}

% --- GESTÃO CLÍNICA ---

\paragraph{Subsistema de Gestão Clínica}

\begin{table}[H]
\centering
\begin{tabularx}{\textwidth}{|l|X|}
\hline
\textbf{Caso de uso} & Submeter Atestado Médico \\
\hline
\textbf{Ator} & Atleta / Encarregado de Educação \\
\hline
\textbf{Descrição} & Atleta ou encarregado submete um atestado médico para validação. \\
\hline
\textbf{Pré-condição} & Utilizador tem sessão iniciada. \\
\hline
\textbf{Pós-condição} & O atestado médico é submetido e fica pendente de validação pela secretaria. \\
\hline
\textbf{Fluxo normal} & 
1. Utilizador acede à secção de atestados médicos. \newline
2. Utilizador carrega o ficheiro do atestado médico. \newline
3. Utilizador indica a data de emissão e validade do atestado. \newline
4. Sistema valida o formato do ficheiro. \newline
5. Sistema regista o atestado com estado ``Pendente''. \newline
6. Sistema notifica a secretaria sobre o novo atestado. \\
\hline
\textbf{Fluxo de exceção 1} & 
[formato de ficheiro inválido] (passo 4) \newline
4.1. Sistema informa sobre os formatos aceites. \newline
4.2. Regressa a 2. \\
\hline
\end{tabularx}
\caption{Caso de uso ``Submeter Atestado Médico''.}
\label{tab:uc_submeter_atestado}
\end{table}

\begin{table}[H]
\centering
\begin{tabularx}{\textwidth}{|l|X|}
\hline
\textbf{Caso de uso} & Validar Atestado Médico \\
\hline
\textbf{Ator} & Secretaria \\
\hline
\textbf{Descrição} & Secretaria valida ou rejeita um atestado médico submetido. \\
\hline
\textbf{Pré-condição} & Secretaria tem sessão iniciada e existem atestados pendentes de validação. \\
\hline
\textbf{Pós-condição} & O atestado é validado ou rejeitado, e o utilizador é notificado. \\
\hline
\textbf{Fluxo normal} & 
1. Secretaria consulta a lista de atestados pendentes. \newline
2. Secretaria seleciona um atestado. \newline
3. Sistema apresenta o documento e os dados do atestado. \newline
4. Secretaria valida o atestado. \newline
5. Sistema atualiza o estado para ``Válido''. \newline
6. Sistema atualiza a situação clínica do atleta. \newline
7. Sistema notifica o atleta/encarregado sobre a validação. \\
\hline
\textbf{Fluxo alternativo 1} & 
[secretaria rejeita o atestado] (passo 4) \newline
4.1. Secretaria indica o motivo da rejeição. \newline
4.2. Sistema atualiza o estado para ``Rejeitado''. \newline
4.3. Sistema notifica o utilizador sobre a rejeição. \\
\hline
\end{tabularx}
\caption{Caso de uso ``Validar Atestado Médico''.}
\label{tab:uc_validar_atestado}
\end{table}

\begin{table}[H]
\centering
\begin{tabularx}{\textwidth}{|l|X|}
\hline
\textbf{Caso de uso} & Registar Lesão \\
\hline
\textbf{Ator} & Treinador \\
\hline
\textbf{Descrição} & Treinador regista uma lesão sofrida por um atleta. \\
\hline
\textbf{Pré-condição} & Treinador tem sessão iniciada e o atleta pertence a uma das suas equipas. \\
\hline
\textbf{Pós-condição} & A lesão é registada no historial clínico do atleta. \\
\hline
\textbf{Fluxo normal} & 
1. Treinador seleciona o atleta. \newline
2. Treinador indica o tipo de lesão, data de ocorrência e descrição. \newline
3. Treinador indica a previsão de recuperação (se conhecida). \newline
4. Sistema regista a lesão. \newline
5. Sistema notifica o encarregado de educação sobre a lesão. \newline
6. Sistema atualiza a disponibilidade do atleta. \\
\hline
\end{tabularx}
\caption{Caso de uso ``Registar Lesão''.}
\label{tab:uc_registar_lesao}
\end{table}

\begin{table}[H]
\centering
\begin{tabularx}{\textwidth}{|l|X|}
\hline
\textbf{Caso de uso} & Reportar Queixa Física \\
\hline
\textbf{Ator} & Atleta / Encarregado de Educação \\
\hline
\textbf{Descrição} & Atleta ou encarregado reporta uma queixa física ao treinador. \\
\hline
\textbf{Pré-condição} & Utilizador tem sessão iniciada. \\
\hline
\textbf{Pós-condição} & A queixa física é registada e o treinador é notificado. \\
\hline
\textbf{Fluxo normal} & 
1. Utilizador acede à secção de queixas físicas. \newline
2. Utilizador descreve a queixa (zona afetada, tipo de dor, intensidade). \newline
3. Utilizador indica se pretende participar no próximo treino/jogo. \newline
4. Sistema regista a queixa. \newline
5. Sistema notifica os treinadores das equipas do atleta. \\
\hline
\end{tabularx}
\caption{Caso de uso ``Reportar Queixa Física''.}
\label{tab:uc_reportar_queixa}
\end{table}

% --- GESTÃO FINANCEIRA ---

\paragraph{Subsistema de Gestão Financeira}

\begin{table}[H]
\centering
\begin{tabularx}{\textwidth}{|l|X|}
\hline
\textbf{Caso de uso} & Gerar Quotas da Época \\
\hline
\textbf{Ator} & Gerência \\
\hline
\textbf{Descrição} & Sistema gera automaticamente as quotas para todos os atletas ativos numa nova época. \\
\hline
\textbf{Pré-condição} & Gerência tem sessão iniciada e foi criada uma nova época desportiva. \\
\hline
\textbf{Pós-condição} & São criadas quotas para todos os atletas ativos, em estado ``Pendente''. \\
\hline
\textbf{Fluxo normal} & 
1. Gerência inicia o processo de geração de quotas para a época. \newline
2. Sistema identifica todos os atletas com inscrição ativa. \newline
3. Sistema calcula o valor da quota para cada atleta (com base no escalão e modalidades). \newline
4. Sistema cria as quotas em estado ``Pendente''. \newline
5. Sistema notifica os encarregados de educação sobre as quotas geradas. \newline
6. Sistema apresenta o resumo da operação. \\
\hline
\end{tabularx}
\caption{Caso de uso ``Gerar Quotas da Época''.}
\label{tab:uc_gerar_quotas}
\end{table}

\begin{table}[H]
\centering
\begin{tabularx}{\textwidth}{|l|X|}
\hline
\textbf{Caso de uso} & Efetuar Pagamento \\
\hline
\textbf{Ator} & Encarregado de Educação \\
\hline
\textbf{Descrição} & Encarregado efetua o pagamento de uma quota pendente. \\
\hline
\textbf{Pré-condição} & Encarregado tem sessão iniciada e existem quotas pendentes associadas aos seus educandos. \\
\hline
\textbf{Pós-condição} & O pagamento é registado e a quota transita para o estado ``Paga''. \\
\hline
\textbf{Fluxo normal} & 
1. Encarregado consulta as quotas pendentes. \newline
2. Sistema apresenta a lista de quotas com valores e prazos. \newline
3. Encarregado seleciona as quotas a pagar. \newline
4. Encarregado escolhe o método de pagamento (MB Way, transferência, numerário). \newline
5. Sistema processa o pagamento. \newline
6. Sistema atualiza o estado das quotas para ``Paga''. \newline
7. Sistema gera o recibo/fatura. \newline
8. Sistema envia confirmação ao encarregado. \\
\hline
\textbf{Fluxo de exceção 1} & 
[pagamento falha] (passo 5) \newline
5.1. Sistema informa sobre a falha no pagamento. \newline
5.2. Sistema mantém as quotas em estado ``Pendente''. \\
\hline
\end{tabularx}
\caption{Caso de uso ``Efetuar Pagamento''.}
\label{tab:uc_efetuar_pagamento}
\end{table}

\begin{table}[H]
\centering
\begin{tabularx}{\textwidth}{|l|X|}
\hline
\textbf{Caso de uso} & Registar Pagamento Manual \\
\hline
\textbf{Ator} & Secretaria \\
\hline
\textbf{Descrição} & Secretaria regista um pagamento efetuado presencialmente ou por transferência. \\
\hline
\textbf{Pré-condição} & Secretaria tem sessão iniciada e existem quotas pendentes. \\
\hline
\textbf{Pós-condição} & O pagamento é registado e a quota transita para o estado ``Paga''. \\
\hline
\textbf{Fluxo normal} & 
1. Secretaria pesquisa o atleta ou encarregado. \newline
2. Sistema apresenta as quotas pendentes. \newline
3. Secretaria seleciona as quotas pagas. \newline
4. Secretaria indica o método de pagamento e comprovativo (se aplicável). \newline
5. Sistema atualiza o estado das quotas para ``Paga''. \newline
6. Sistema gera o recibo/fatura. \newline
7. Sistema notifica o encarregado sobre o pagamento registado. \\
\hline
\end{tabularx}
\caption{Caso de uso ``Registar Pagamento Manual''.}
\label{tab:uc_registar_pagamento}
\end{table}

\begin{table}[H]
\centering
\begin{tabularx}{\textwidth}{|l|X|}
\hline
\textbf{Caso de uso} & Gerar Relatório SAF-T \\
\hline
\textbf{Ator} & Gerência \\
\hline
\textbf{Descrição} & Sistema gera o ficheiro SAF-T para submissão à Autoridade Tributária. \\
\hline
\textbf{Pré-condição} & Gerência tem sessão iniciada e existem transações financeiras registadas. \\
\hline
\textbf{Pós-condição} & É gerado um ficheiro SAF-T válido com todas as transações do período selecionado. \\
\hline
\textbf{Fluxo normal} & 
1. Gerência indica o período para o relatório (mês/ano). \newline
2. Sistema recolhe todas as transações do período. \newline
3. Sistema valida a conformidade dos dados. \newline
4. Sistema gera o ficheiro SAF-T em formato XML. \newline
5. Sistema disponibiliza o ficheiro para download. \\
\hline
\textbf{Fluxo de exceção 1} & 
[dados incompletos ou inválidos] (passo 3) \newline
3.1. Sistema apresenta lista de erros/inconsistências. \newline
3.2. Sistema cancela a geração do ficheiro. \\
\hline
\end{tabularx}
\caption{Caso de uso ``Gerar Relatório SAF-T''.}
\label{tab:uc_gerar_saft}
\end{table}

\begin{table}[H]
\centering
\begin{tabularx}{\textwidth}{|l|X|}
\hline
\textbf{Caso de uso} & Consultar Situação Financeira \\
\hline
\textbf{Ator} & Encarregado de Educação / Atleta \\
\hline
\textbf{Descrição} & Utilizador consulta o estado das suas quotas e histórico de pagamentos. \\
\hline
\textbf{Pré-condição} & Utilizador tem sessão iniciada. \\
\hline
\textbf{Pós-condição} & São apresentadas as quotas e pagamentos associados ao utilizador. \\
\hline
\textbf{Fluxo normal} & 
1. Utilizador acede à secção financeira. \newline
2. Sistema apresenta o resumo da situação (quotas pendentes, em atraso, pagas). \newline
3. Sistema apresenta o histórico de pagamentos e faturas. \\
\hline
\end{tabularx}
\caption{Caso de uso ``Consultar Situação Financeira''.}
\label{tab:uc_consultar_financeira}
\end{table}

% ============================================
% DIAGRAMAS DE ATIVIDADE
% ============================================

\subsubsection{Diagramas de atividade}

Para a modelação do software, foram também construídos diagramas de atividade para as operações 
com fluxos de controlo mais complexos. Estes diagramas descrevem operações da camada de negócio, 
permitindo uma representação organizada do controlo de fluxo que facilita a implementação.

\paragraph{Diagrama de atividade: Registo de utilizador}

O diagrama seguinte descreve a resposta do sistema ao pedido de registo de um novo utilizador:

\begin{figure}[H]
    \centering
    \includegraphics[width=0.85\textwidth]{imagens/atividade_registo.png}
    \caption{Diagrama de atividade relativo ao registo de um utilizador.}
    \label{fig:atividade_registo}
\end{figure}

A atividade inicia-se com a receção dos dados do utilizador. O sistema verifica se já existe um 
utilizador com o email providenciado. Em caso afirmativo, o sistema informa o erro e termina. 
Caso contrário, calcula a hash da palavra-passe, regista os dados com estado ``Pendente'', 
notifica a secretaria e redireciona o utilizador para a página de confirmação.

\paragraph{Diagrama de atividade: Processamento de pagamento}

O diagrama seguinte descreve o fluxo de processamento de um pagamento de quotas:

\begin{figure}[H]
    \centering
    \includegraphics[width=0.9\textwidth]{imagens/atividade_pagamento.png}
    \caption{Diagrama de atividade relativo ao processamento de pagamento.}
    \label{fig:atividade_pagamento}
\end{figure}

Após receber os dados do pagamento, o sistema valida as quotas selecionadas em paralelo com a 
verificação do método de pagamento. Se ambas as validações forem bem-sucedidas, o sistema processa 
o pagamento junto do gateway. Em caso de sucesso, atualiza o estado das quotas, gera a fatura e 
envia a confirmação. Em caso de falha, informa o utilizador do erro.

\paragraph{Diagrama de atividade: Publicação de convocatória}

O diagrama seguinte descreve o processo de publicação de uma convocatória e envio de notificações:

\begin{figure}[H]
    \centering
    \includegraphics[width=0.85\textwidth]{imagens/atividade_convocatoria.png}
    \caption{Diagrama de atividade relativo à publicação de uma convocatória.}
    \label{fig:atividade_convocatoria}
\end{figure}

A atividade inicia-se com a receção do pedido de publicação. O sistema valida que a convocatória 
está em estado ``Rascunho'' e que tem pelo menos um atleta convocado. Se válido, atualiza o estado 
para ``Publicada'' e, em paralelo, envia notificações para cada atleta convocado e respetivo 
encarregado de educação. A estrutura cíclica itera por todos os atletas da convocatória.

\paragraph{Diagrama de atividade: Verificação de quotas em atraso}

O diagrama seguinte descreve o processo automático de verificação de quotas em atraso:

\begin{figure}[H]
    \centering
    \includegraphics[width=0.85\textwidth]{imagens/atividade_quotas_atraso.png}
    \caption{Diagrama de atividade relativo à verificação de quotas em atraso.}
    \label{fig:atividade_quotas_atraso}
\end{figure}

Este processo é executado periodicamente pelo sistema. Para cada quota em estado ``Pendente'', 
o sistema verifica se o prazo de pagamento expirou. Em caso afirmativo, atualiza o estado para 
``Em Atraso'' e envia uma notificação ao encarregado de educação. Se o atraso ultrapassar um 
limite configurável, o sistema pode sinalizar o atleta para suspensão de atividades.

